The increasing reliance on lithium-ion batteries, particularly in the automotive industry, necessitates the insurance of safety and performance of these energy storage systems. Lithium-ion pouch cells, widely used in electric vehicles, present specific challenges when subjected to mechanical loading. Large deformations of the mechanically loaded cells lead to potential safety hazards such as fire risks. Thus, modeling the mechanical behavior of these cells under mechanical load is crucial to predicting their response and mitigating potential risks.

Previous studies have used isotropic pressure-dependent plasticity models to describe the mechanical behavior of lithium-ion cells. However, these models often rely on simplified hardening laws that may not fully capture the complexity of the cell’s response to deformation. The hardening models, typically expressed through two-stage analytical models, may limit the ability to generalize the response under different loading conditions, particularly for complex geometries or diverse indentation scenarios. This thesis aims to overcome these limitations by leveraging deep learning techniques to predict the hardening response of lithium-ion pouch cells more accurately.

The primary objective of this research is to develop a neural network-based hardening law that can be integrated into finite element simulations, enabling more precise prediction of yield stress from plastic strain. By utilizing a fully connected neural network (FCNN), contrary to traditional hardening models and this approach relies on the neural network's ability to learn nonlinear relationships from data. The proposed model is designed to enhance computational efficiency while providing a more accurate representation of the material behavior, especially in complex indentation scenarios.

This thesis outlines the development of using a deep learning framework for calibrating the hardening law based on experimental data from indentation tests. The neural network is trained to predict the deformation resistance as a function of the equivalent plastic strain, with the calibration process relying on a fast simulation framework. The approach is further extended to investigate the performance of the model for various loading conditions and geometries.

In addition to the primary objective of improving force-displacement matching, this work explores the limitations of current modeling approaches, including convergence issues and the oscillating loss curves observed in previous studies. By incorporating machine learning into the calibration process, this research aims to establish a more robust and general framework for simulating the mechanical behavior of lithium-ion pouch cells. Thus, contributing to advancements in material science and battery technology \citep{Tancogne-Ind}.
